\begin{flushright}
\huge{\textbf{Abstract}}
\end{flushright}
\addcontentsline{toc}{section}{Abstract}

Quantitative equity trading models frequently fail out-of-sample due to survivorship bias, multicollinear feature representations, and unaccounted transaction frictions. This thesis presents the design, empirical evaluation, and live deployment of \textbf{XGForge}, an end-to-end quantitative equity trading and portfolio management framework implemented on the Indian equity market (NSE NIFTY 200). 

We design an object-oriented backtesting and execution infrastructure that decouples data ingestion, feature generation, forward target formulation, statistical estimation, portfolio selection, and risk management into swappable modular contracts. To eliminate survivorship bias, we engineer a dynamic point-in-time constituent mapping pipeline tracking 260 historical symbols across index reconstitutions from 2000 to 2026. 

For alpha discovery, we implement a two-stage feature selection pipeline. An initial candidate universe of 60 technical indicators across six factor families is ranked by tree-based split gain and pruned against pairwise collinearity. A cumulative gain scree plot confirms that a parsimonious subset of 16 high-gain indicators captures 86.4\% of total predictive signal. Using a chronological training split, we benchmark non-parametric gradient boosted decision trees (XGBoost) against linear estimators (Ridge, Lasso) and ensemble alternatives (Random Forest, LightGBM) on a 63-day forward return target. XGBoost achieves superior out-of-sample forecast accuracy ($R^2_{\text{OOS}} = 0.94\%$) and lowest RMSE (0.1973), while LightGBM and Random Forest achieve higher Rank Information Coefficients (0.0425 and 0.0402). Interpretability analysis via TreeSHAP confirms that medium-term momentum and negative realized volatility drive positive expected alpha.

Across an automated search harness executing 9,182 backtest simulations (Grid, Explicit, and One-Factor-At-A-Time modes), XGForge ranked first by the composite dual-hurdle score (coming closest to clearing the dual-hurdle criteria where no configuration cleared both). In an out-of-sample event-driven simulation (2019--2026) audited under realistic transaction frictions, dynamic bid-ask slippage (\textsf{Rs.~}3.64 Lakhs on \textsf{Rs.~}12.03 Crores turnover), and statutory capital gains tax (\textsf{Rs.~}4.21 Lakhs) across 5,547 fills, the strategy generates a net compound annual growth rate (CAGR) of 17.08\% (pre-tax 19.04\%) versus 11.19\% for the NIFTY 200 benchmark. The strategy improves the annualized Sharpe ratio from 0.712 to 1.059 and cuts maximum drawdown from -38.22\% to -20.74\%, preserving capital through the March 2020 liquidity shock. Finally, we deployed a production variant of XGForge (same XGBoost regressor, 16 multi-factor features, and expanding walk-forward training; configured for top-10 stock selection and 95\% capital deployment) into OpenAlgo’s sandbox operating on live Angel One prices; paper-trading evaluation from November 2 to November 15, 2026 EOD recorded an absolute portfolio return of +8.56\% (+\textsf{Rs.~}8.13 Lakhs on \textsf{Rs.~}94.96 Lakhs deployed) across 10 trading sessions with all 10 positions profitable.

\vspace{0.5cm}
\noindent\textbf{Keywords:} Quantitative Finance, XGForge, Gradient Boosted Decision Trees, XGBoost, Survivorship Bias, TreeSHAP, Market Frictions, Event-Driven Backtesting, OpenAlgo.
\clearpage
